% Chapter 9: Energy: Accounting for Change (complete sample chapter)
\chapter{Energy: Accounting for Change}

\step{A Counterintuitive Opening}

\begin{mainline}
An American physics professor performed a famous classroom demonstration. He suspended a steel ball weighing more than ten kilograms from the ceiling on a thick steel cable, making a giant pendulum. Backing against the wall, he pulled the ball to his nose, touching the skin, and released it, without pushing. It swept across the room and rushed back toward his face. The class held its breath. A few centimeters short of his nose it stopped and swung away again. The professor never flinched. He said he trusted physics, so he did not dodge.

The part deserving sleepless thought is the ball's discipline, not the professor's courage. Why can it \emph{never} return higher than its starting point? It knows nothing of noses; it is dead metal, yet obediently stops slightly lower every time. Conversely, give it a secret push at release and it swings higher, genuinely hitting the nose. Where does that extra height come from, and how does the ball remember its original height so faithfully?

You can watch a version anytime: a park swing with no pushing or pumping only loses amplitude until it stops at the bottom. No swing anywhere spontaneously climbs higher. Water flows downhill and pendulums settle low, as though nature keeps an invisible account: something changes hands, but its total only decreases, never increases, and nobody gets extra for nothing. What is recorded? This chapter opens the ledger.
\end{mainline}

\step{Make a Guess First}

As usual, bet first and read second. Write your answers.

Question one. Identical balls start from rest at one height, one down a smooth straight slope, the other down a smooth winding track that dips and rises along the way. Ignore friction and air resistance. Which reaches the bottom faster in speed? (a) The straight-slope ball; (b) the curved-track ball; (c) both have the same speed.

Question two. A car traveling at \(40\) kilometers per hour brakes hard and slides \(10\) meters to rest. On the same road at \(80\) kilometers per hour, what braking distance should it need? (a) \(20\) meters; (b) \(40\) meters; (c) less than \(40\).

Question three. An inventor builds a vertical wheel with swinging weights on its rim, claiming they always extend farther out on the right and retract on the left. The right is therefore always heavier, so the wheel turns forever while driving a generator. Can it work? If yes, why has nobody generated electricity with it for centuries? If no, where does it fail?

Record your judgments first. We will reconcile them at the end.

\step{Hands-on Experiment}

\begin{experiment}[A Key Pendulum: Does It Remember Height?]
\textbf{Materials}: thin string about a meter long, an earphone cable or sewing thread will do; a key or nut; and a pen.

\textbf{Step one}: tie the key to one end, mark the other end with the pen, and grip it securely so the key hangs vertically. This is a pendulum.

\textbf{Step two}: pull the key beside your nose with the string taut, then \emph{release gently}, never pushing. It swings away and returns. Watch its highest return point. Dare you use your nose as a reference? Provided you really did not push, it is safe, the professor's demonstration at home.

\textbf{Step three}: try another release height. Wherever you start, the key almost returns to the \emph{same} height, always slightly lower. Each successive swing reaches lower until it stops at the bottom.

\textbf{Step four}: hold a pen across the string halfway down, intercepting it after the bob passes its lowest point and abruptly shortening the swinging radius. Remarkably, the key still climbs almost to the original \emph{height} on the other side, despite following a different arc.

\textbf{Record}: what quantity does it remember? Velocity? Bottom speed changes each time. Path? The pen changes it instantly. Height, however, is remembered clearly. We will ask what account the world is keeping.
\end{experiment}

\begin{experiment}[Rubber-Band Launcher: Twice the Stretch, How Much Energy?]
\textbf{Materials}: a thick rubber band, a small paper rod made by tightly rolling a sticky note, a ruler, and a book as a launching platform.

\textbf{Step one}: fix one band end under the book's spine and hook the other around the rod. Pull horizontally along the table, measuring \(2\) centimeters of extension, release, and measure landing distance from the table edge. Repeat three times and average.

\textbf{Step two}: extend it \(4\) centimeters, exactly double, and average three launches again. How many times farther will it land? Most intuitions predict roughly twice.

\textbf{Step three}: measurements usually show about three times or more. Landing distance depends on release speed, and doubling speed doubles range. A threefold range means roughly doubled release speed, with fourfold \emph{kinetic energy} carried by the rod.

\textbf{Record}: extension only doubled, yet energy more than quadrupled. How is energy stored in a band? The Mathematical Translator will draw this square relationship as a triangle.
\end{experiment}

\step{The Old Model Runs into Trouble}

Through the previous chapter, our tools were forces: free-body diagrams, Newton's laws, and specific force types. They answer present-instant questions superbly: what force now, what acceleration now? But many everyday questions concern the accumulated result over a whole process. Force language shows three shortcomings there.

Difficulty one: the same force gives unexpectedly different results. Push one cart with the same force for one second or two and velocity differs, understandably growing roughly with time. Measure how far it can then coast and something seems wrong: twice the speed gives roughly four times the coasting distance, not twice. If intuition says force gave the cart a store of impetus that runs out, that impetus refuses to reconcile when counted by speed versus distance. The old language has no quantity unifying how long you pushed and how far it coasted.

Difficulty two: bullets through wood. A bullet can just penetrate one board, ending with zero speed. Can two identical bullets at that speed penetrate two stacked boards? Yes, but after the first board the speed \emph{does not} halve. Intuition says each board subtracts the same speed; experiment and theory say it subtracts something else, leaving speed about seventy percent of the original after that subtraction. What is this something? Force language has no place for it.

Difficulty three: remembered height. Change a key pendulum's height, radius, or intercept its string, and it always recognizes its original horizontal level. Force changes at every point as tension direction changes. Change the track and every force detail changes, yet step through Newton's laws and the final height stays the same. It resembles driving winding mountain roads to a neighboring city with a fuel reading determined solely by starting and finishing points, regardless of route. The route, the forces, varies endlessly while one quantity stays fixed. That quantity cannot be force itself; it sees a larger picture.

The three difficulties share a shape: forces excel at instant-by-instant tracking but less at \emph{summarizing a whole process}. There are two ways to summarize. Accumulate force over \emph{time} and obtain impulse, the next chapter's collision protagonist. Accumulate it over \emph{distance} and obtain this chapter's work. Newton's second law favors neither, but nature accepts both, so we need two ledgers. Start with distance accounting: ask not every instant's push and pull but what moves from place to place or changes form over the entire process while the total remains exact. Physicists searched nearly two hundred years for this quantity, giving it odd names such as living force along the way. Today we call it energy.

\step{Inventing New Concepts}

\begin{mainline}
\textbf{First concept: work, a force's transfer slip.} To summarize a process, define how much a force accomplishes along a distance. Push a door along its motion and it opens; carry a full bucket horizontally and, however sore your arm, the bucket gains no height. These examples suggest counting only force \emph{along displacement}. Define \emph{work} as force magnitude times the object's displacement along that force. When carrying the bucket, upward force and horizontal displacement are perpendicular, so that force does zero work. Soreness is physiological consumption inside muscles, not a transfer to the bucket. Work resembles a transfer slip: force is the act of paying, displacement the payment distance, their product the transferred amount. It is unlike a wallet: an object cannot contain a hundred joules of work, any more than you can carry a hundred dollars of bank transfer. Transfer happens only during a process.

\textbf{Second concept: energy, the account balance.} Transfer slips need accounts. Define \emph{energy as a measure of a system's capacity to do work on its surroundings}. A raised hammer drives a pile, a wound spring runs a watch, a flying bullet pierces wood: each has a balance. The definition's subtlety is stating what energy can do, not what stuff it is. Energy is no tangible substance but a \emph{bookkeeping quantity} invented by physicists. Do not underestimate one: we will see the universe supervises this ledger most strictly.

\textbf{Third concept: kinetic energy, motion's balance.} The bullet puzzle now has an answer. A moving object's balance is proportional to \emph{speed squared}, not speed. Call it kinetic energy, \(\frac{1}{2}mv^2\). Double speed and kinetic energy quadruples, so braking distance quadruples. Stretch the rubber band twice as far and the launched paper receives fourfold kinetic energy, more precisely fourfold supplied energy, carrying it much farther. This gives question two's counterintuitive answer (b).

\textbf{Fourth concept: potential energy, a stored balance.} Where is the energy of a raised hammer before it falls? Neither in its molecules nor in the air, but in the \emph{positional relationship} of hammer plus Earth. Energy stored by position is potential energy. Near ground, gravitational potential energy is weight times height. A stretched spring or band stores elastic potential energy roughly proportional to deformation squared. It resembles banked money, redeemable with interest, actually to the exact penny, as kinetic energy. It differs from cash: you cannot discuss it outside the hammer--Earth or spring--hand system. Saying an object has gravitational potential energy is convenient shorthand; the actual account belongs to the system.

\textbf{Fifth concept: conservative and nonconservative forces, reversible and irreversible transfers.} Gravity and springs are obliging: lift an object and return it, and their positive and negative work cancel. Every transferred penny can return. Such forces are \emph{conservative}, allowing potential energy definitions. Friction does the opposite: push a box out and back and it charges a fee both ways, reducing mechanical energy. It is \emph{nonconservative}. Where did its takings go? They did not vanish. Feel a brake disc after hard braking or rub cold hands together: they warm. Energy changes form into ceaseless molecular jostling, the account called internal energy, central to thermal physics. Thus \emph{mechanical energy can decrease; total energy never does}. Friction destroys only a useful form, not energy itself.

\textbf{Sixth concept: power, the transfer rate.} Walking and running up three flights transfer the same weight-times-height amount, yet leave you breathing differently. The distinction is \emph{rate}: work per unit time is power. A weightlifter raises a bar in one or two seconds with impressive power; an elderly woman taking ten minutes to raise the same load the same height does identical work at far lower power. Your bill's kilowatt-hours measure energy; an appliance's watts measure power. Balance and transfer rate are different.
\end{mainline}

All six concepts are ready. Return to the invisible ledger. At the professor's nose the ball has gravitational potential energy; at the bottom it has become kinetic; at the opposite top kinetic converts fully back to potential, explaining the remembered height. The slight loss each time goes through air drag and pivot friction to internal energy, so returns can be lower, never higher. Transfers are allowed; issuing extra money is not.

Use this language and life becomes a sequence of transfers. Drawing a bow converts muscular chemical energy to elastic potential, mostly becoming the arrow's kinetic energy on release. Powerful ancient bows made this conversion fast and fierce enough for arrows to exceed high-speed train speeds. A trampoline stores falling kinetic energy elastically and returns it upward; each bounce's losses become fabric and air internal energy, so without pumping you rise less. Pendulums, swings, diving boards, and pole-vault poles repeat the same drama: potential and kinetic exchange while friction takes commission. Conversely, an escalator raises you at constant velocity, unchanged kinetic energy but increasing gravitational potential. Its motor transfers steadily at constant power. Lose electrical supply and your chemical energy must pay, making your legs feel heavy at once.

\step{The Model in One Sentence}

\begin{oneline}
Energy is the universe's universal currency: never created or destroyed, only transferred between kinetic, potential, internal, and other accounts. Work and heat are transfer methods, power the transfer rate, and mechanical energy is conserved only when conservative forces handle the account.
\end{oneline}

Every word corresponds to a preceding definition. The currency analogy matches conservation, interchangeability, storage, and circulation. It fails because energy has no inflation, the universe never prints money, and no bank lends from nothing. That is why perpetual-motion machines cannot work.

\step{The Mathematical Translator}

Translate every plain-language statement. Before each formula, ask its question; afterward, immediately test special cases.

\textbf{Work: force accumulated along a path.} How much does a constant force transfer along straight-line displacement? If \(\mathbf F\) and \(\mathbf s\) make angle \(\theta\), only the component \(F\cos\theta\) along displacement counts, giving
\[
W = Fs\cos\theta .
\]
Check: \(\theta=0\), pushing along motion, gives \(W=Fs\), the full amount. At \(\theta=90^{\circ}\), carrying a bucket horizontally or string tension in circular motion, \(\cos\theta=0\) and \(W=0\): inward force turns without transferring. This explains constant speed in circular motion with no other tangential force. At \(\theta=180^{\circ}\), sliding friction gives \(W=-Fs\), the minus sign recording withdrawal from mechanical energy. Work is scalar but signed: positive and negative mean transfer in or out.

\begin{center}
\begin{tikzpicture}[scale=1.0,>={Stealth[length=3mm]}]
  \draw[fill=gray!15] (0,0) rectangle (1.2,0.8);
  \draw[->,very thick] (0.6,0.4) -- (3.1,1.84) node[above left]{$\mathbf F$};
  \draw[->,very thick] (0.6,0.4) -- (3.6,0.4) node[midway,below]{$\mathbf s$};
  \draw[dashed] (2.6,0.4) arc[start angle=0,end angle=30,radius=2.0];
  \node at (2.45,0.75) {$\theta$};
  \draw[dashed] (3.1,1.84) -- (3.1,0.4);
  \draw[->,thick] (0.6,0.4) -- (3.1,0.4) node[midway,above=2pt]{\small $F\cos\theta$ effective component};
\end{tikzpicture}
\end{center}

\textbf{Kinetic energy: why speed squared?} How much balance does motion carry? Existing knowledge derives it. Let constant \(F\) push mass \(m\) from rest a distance \(s\) along the force. Newton gives \(a=F/m\); uniform-acceleration kinematics gives \(v^2=2as\). Eliminate \(a\):
\[
Fs=\frac{1}{2}mv^2 .
\]
The left is the external transfer in, so the right must be the new balance. Thus kinetic energy is
\[
E_{\mathrm k}=\frac{1}{2}mv^2 .
\]
Check: \(v=0\) gives zero, reasonably. Double \(v\) and energy quadruples: braking at \(80\) kilometers per hour takes about four times the distance at \(40\), not twice. Hence highway following margins must scale with speed squared; question two's intuition is overturned.

\begin{mathbridge}
\textbf{Variable-force work: area and integration.} For force varying with position, divide the path into tiny intervals with nearly constant force. Each work increment is \(\Delta W\approx F(x)\Delta x\); total work is area under the \(F\)-\(x\) graph:
\[
W=\int_{x_1}^{x_2}F(x)\,\mathrm dx .
\]
The same infinite subdivision and summation for variable acceleration proves the general \emph{work--energy theorem}: total work by net external force equals kinetic-energy change,
\[
W_{\text{net}}=\Delta E_{\mathrm k}=\frac{1}{2}mv_2^2-\frac{1}{2}mv_1^2 .
\]
It applies to constant or variable forces and straight or curved paths. That is why energy accounting saves effort over instant-by-instant force analysis: ask endpoints, not process.
\end{mathbridge}

\textbf{Gravitational potential energy: weight times height.} Raising an object by \(h\) stores how much? Work against gravity is \(mgh\), all deposited, so
\[
E_{\mathrm p}=mgh ,
\]
with \(h\) measured from an arbitrarily chosen zero height. Check: at \(h=0\), potential is zero. The origin is bookkeeping convention; only height \emph{differences} matter. Negative \(h\) is legal too: a basement object has a negative balance, a debt, yet height differences remain usable. Hydroelectric plants exploit exactly this: water mass times drop times \(g\) is the reservoir's balance.

\textbf{Elastic potential energy: the half in the area.} How much is stored by stretching a spring \(x\) from its natural length? Hooke says pull grows linearly, \(F=kx\). It varies from \(0\) to \(kx\), with mean \(\frac{1}{2}kx\); multiply by distance \(x\):
\[
E_{\mathrm p}=\frac{1}{2}kx^2 .
\]
The half is the triangle's area on the force--displacement graph, not an inserted convenience. Check: \(x=0\) gives zero; \(\pm x\) gives identical values, equal storage in compression and extension. The launcher's twice-the-stretch, four-times-the-energy puzzle lies in this square.

\begin{center}
\begin{tikzpicture}[scale=0.9,>={Stealth[length=2.5mm]}]
  \draw[->] (0,0) -- (4.6,0) node[right]{$x$};
  \draw[->] (0,0) -- (0,3.4) node[above]{$F$};
  \draw[very thick] (0,0) -- (3.6,2.7) node[right]{$F=kx$};
  \draw[dashed] (3.6,0) node[below]{$x$} -- (3.6,2.7) -- (0,2.7) node[left]{$kx$};
  \fill[pattern=north east lines,pattern color=blue!60] (0,0) -- (3.6,2.7) -- (3.6,0) -- cycle;
  \node at (2.55,0.85) {$E_{\mathrm p}=\dfrac12 kx^2$};
\end{tikzpicture}
\end{center}

\textbf{Mechanical energy conservation: only conservative forces doing work.} When is kinetic plus potential constant? Account for a process with only conservative-force work: gravity adds exactly as much kinetic energy as it removes potential. Thus
\[
E_{\mathrm k}+E_{\mathrm p}=\text{constant}\qquad(\text{only conservative forces do work}) .
\]
Check extremes: at a pendulum's top, \(v=0\) and all energy is potential; at the bottom, potential is least and kinetic greatest. Between them one rises as the other falls, total unchanged. Add friction or drag and the right side no longer stays constant: mechanical losses accrue. They have not evaporated but entered internal energy. The complete law says \emph{include internal energy and the total always balances}.

\textbf{Power: transfer rate.} How quickly does work happen over time?
\[
P=\frac{W}{t},\qquad\text{for constant force along velocity, also }P=Fv .
\]
Check two extremes of \(P=Fv\): however hard you push an immovable wall, \(v=0\) gives zero power and no joule reaches its account. A crane lifting a container \emph{uniformly} has zero net force but nonzero cable tension, transferring energy into the container--Earth system at \(P=Fv\). This also explains a driving fact: twice the speed requires twice the power for unchanged tractive force, while power against air drag grows still faster, the first reason highway speeds burn fuel.

\begin{mathbridge}
\textbf{Power as a derivative.} More precisely, power is work's rate of change with time:
\[
P=\frac{\mathrm dW}{\mathrm dt}=\mathbf F\cdot\frac{\mathrm d\mathbf s}{\mathrm dt}=\mathbf F\cdot\mathbf v ,
\]
The dot product automatically selects the force component along motion, gathering every earlier angle case into one symbol: perpendicular means no work, opposite means negative work.
\end{mathbridge}

\textbf{The electricity bill's ledger.} Electricity charges are energy accounting's most down-to-earth application. A \(1\)-kilowatt appliance running \(1\) hour consumes \(1\) kilowatt-hour, commonly one unit of electricity:
\[
1\ \text{kWh}=1000\ \text{W}\times3600\ \text{s}=3.6\times10^{6}\ \text{J} .
\]
How much is that? Used entirely against gravity, it raises one tonne of water about \(360\) meters. A \(1500\)-watt kettle boiling water for about \(5\) minutes uses roughly \(0.12\) kilowatt-hours. An inverter air conditioner's few daily units in cooling season represent a compressor spending hundreds of watts all night moving heat outdoors. It pays a transport fee, not heat's purchase price, the deeper reason air conditioning saves electricity compared with resistance heating, explained in thermal physics.

\textbf{An estimate with real data.} A bowl of rice holds about \(1\times10^{6}\) joules of chemical energy, roughly \(250\) kilocalories. At hypothetical hundred-percent conversion to stair-climbing work, how high could a \(60\)-kilogram person climb? With \(mgh=10^{6}\) joules and \(g\approx10\) meters per second squared,
\[
h=\frac{10^{6}}{60\times10}\approx 1700\ \text{m} ,
\]
equivalent to twice Mount Tai's elevation gain. Real bodies are far from fully efficient; most energy ends as body heat and internal energy involved in sweat evaporation. But the order is honest: food's chemical account is much larger than our feeling of one bowl suggests. For power, the same person taking \(15\) seconds to climb \(10\) meters produces about \(60\times10\times10/15=400\) watts mechanically, briefly more than half a microwave's rated power. Over a full day, human average power is only around a hundred watts. A real horse sustains over seven hundred, the origin of horsepower.

\textbf{An engineering use: turning braking into charging.} Conventional brakes turn the car's kinetic energy into internal energy through pad friction, wasting it as heat. Electric cars and high-speed trains reverse the approach with regenerative braking: wheels drive the motor as a generator, returning kinetic energy to battery or grid. A \(1.5\)-tonne car at \(100\) kilometers per hour, about \(28\) meters per second, has
\[
E_{\mathrm k}=\frac12\times1500\times28^{2}\approx 5.9\times10^{5}\ \text{J}\approx 0.16\ \text{kWh} .
\]
One stop seems modest, but dozens or hundreds in a city commute recover enough for substantial extra range. A train braking into a station recovers enough to run lighting and air conditioning for a while. This creates no energy; it avoids burning the account away. Conservation survives and waste falls.

\step{The Model's Limits}

Every conservation law is a contract with conditions. Energy accounting has five main boundaries.

First, \textbf{mechanical energy conservation is conditional}. Only conservative forces such as gravity and elasticity doing work leave kinetic plus potential constant. With friction, collision deformation, human pushes, or an engine, rewrite it as mechanical-energy change equals nonconservative work. Designing a whole roller coaster as mechanically conservative underestimates the first hill's height and courts disaster; engineers must budget every segment's friction.

Second, \textbf{total energy conservation requires a closed system}. Before saying conservation, locate the ledger's boundary. Enclose only a cart and friction transfers to ground and air, so the cart loses energy. Enclose ground, air, sound, and heat too and the total balances. Here dies the first kind of perpetual-motion machine: claiming a machine both exports electricity and keeps running claims growth in a closed balance. Patent offices received thousands over two hundred years, none verified experimentally. France's Academy of Sciences announced in 1775 that it would no longer examine perpetual-motion patents, not from arrogance but because the ledger had never lost a penny. A subtler second kind respects energy conservation while seeking to convert all the sea's internal energy unconditionally into work. Entropy in thermal physics will sentence that one; keep the suspense for now.

Third, \textbf{potential energy's zero and ownership}. Its zero is arbitrary and only differences matter; it belongs to a system, not one object. Saying a book has so many joules of gravitational potential implicitly selects both a zero, perhaps the tabletop, and the book--Earth system. Without them the sentence is incomplete.

Fourth, \textbf{correct everyday expressions}. A dead battery has not used up energy; it has depleted chemical energy convertible into electricity, which has become light, sound, and the phone's warmth without losing a penny. A machine consumes energy is similar shorthand: energy transfers, often into unwanted forms, rather than vanishing. Meanwhile energetic or full of vitality uses psychological metaphor, unrelated to the measurable, transferable, conserved physical quantity; neither endorses the other. Beware energy as a mysterious fluid too. It is not stuff in an object but our unified measure of change. Its reality lies in an always-balanced account, not the ability to scoop out a spoonful.

Fifth, \textbf{classical accounting's domain}. Every formula assumes speeds far below light speed. Near it, \(\frac{1}{2}mv^2\) fails and energy--mass accounts require relativity, previewed in the challenge box. Microscopic energy remains conserved, but accounts become discrete energy levels, quantum physics' story. Also, conservation is more than an empirical rule without a known counterexample. In the early twentieth century, Noether proved that laws unchanged by time translation, experiments agreeing today and tomorrow, imply a conserved quantity: energy. Energy conservation is the receipt for time's uniformity.

\begin{challenge}
\textbf{The deeper origin of kinetic energy's square law.} Relativity writes total energy as \(E=\gamma mc^{2}\), with light speed \(c\) and \(\gamma=1/\sqrt{1-v^{2}/c^{2}}\), slightly above \(1\) at low speed. Expand \(\gamma\) in powers of \(v^{2}/c^{2}\). The first term is constant \(mc^{2}\), an enormous balance already present at rest; the second is exactly \(\frac{1}{2}mv^{2}\), our kinetic energy as relativity's first low-speed correction. Higher terms are no fine print in accelerators: at \(99.9999991\%\) of light speed, a proton has thousands of times the kinetic energy predicted classically. The ledger framework persists; the exchange rates change.
\end{challenge}

\step{Explaining the Opening Phenomenon}

Return to the classroom steel ball and close the chapter's final loop.

At release beside the nose, speed is zero and all energy is gravitational potential \(mgh\). During descent gravity does positive work, transferring potential into kinetic. At the bottom height is least and speed greatest, with almost everything kinetic. As it rises, transfers reverse. In an ideal closed world without drag or pivot friction, it returns exactly to \(h\), not a millimeter more or less. \emph{Its memory of height is mechanical-energy conservation doing the bookkeeping}. At every instant kinetic plus potential equals initial \(mgh\). At the returning highest point, velocity reverses and kinetic energy must be zero, requiring full restoration of potential and hence height.

Real drag and friction take a small commission into ball, air, and bearing internal energy, so every return is slightly lower. Only lower, never higher, lets the professor pledge his face. A pendulum spontaneously rising would mean a ledger printing money and require rebuilding physics from its foundations.

Why can a swing rise higher, then? The swinger is not a closed ledger. Squatting at the top and standing at the bottom, kicking and shifting the body, converts muscular chemical energy into positive work on swing plus person at sensitive phases of motion. Deposits exceeding frictional fees accumulate greater amplitude. Stop moving and only fees remain, inevitably lowering the swing. Transferring at the right time is a whole subject, resonance, treated in oscillations. For now, changing amplitude is never a miracle, only someone continually topping up.

Question one is settled too. Without friction, the straight and curved tracks have the same height difference and gravitational transfer. From \(\frac{1}{2}mv^2=mgh\), both speeds are \emph{identical}, regardless of track shape: answer (c). Conservative forces recognize endpoints, not routes. Add friction back and the longer, more strongly pressing curved track instead gives a slower arrival. Intuition stumbles again over process not mattering. The pen-intercepted key likewise returns to its original height: radius and track changed, but bottom speed and height difference did not. The ledger checks height differences only.

\step{Thought Experiments and Challenges}

\begin{thought}[Drop a Stone from the Moon]
Imagine standing at lunar orbital distance, about \(3.8\times10^{8}\) meters from Earth's center, and gently \emph{dropping} a stone, without pushing, initially ignoring atmosphere. It accelerates toward Earth; what speed just before impact? Account for it without tracking each segment: gravitational potential decreases from far to near and becomes kinetic. The complete gravitational potential formula, Chapter 8's gift, gives about \(11\) kilometers per second, exactly the second cosmic velocity. No coincidence: an object thrown vertically from ground at \(11\) kilometers per second just escapes to infinity, and falling in reverses flying out in time, so the accounts must agree. Go further: in the atmosphere that huge kinetic balance almost entirely becomes internal energy. That makes meteors burn and demands spacecraft heat shields. A gentle release becomes a fireball on arrival. The energy ledger has no mercy for romantic imagination.
\end{thought}

\begin{puzzles}
\item \textbf{A roller-coaster loop.} A loop has radius \(R\). Release a car from rest on a smooth high slope. How far above the loop's top must the release point be to pass it safely without losing contact? (Hint: nonzero speed is insufficient; the track must supply the centripetal role. At the critical point normal force is zero and gravity alone provides required acceleration, \(v^{2}/R\geq g\). Combine this dynamic condition with conservation to get an extra \(R/2\). Real coasters need a safety margin and start substantially above this threshold.)
\item \textbf{Two arrows.} From the same position at the same speed, launch one arrow vertically upward and another horizontally, ignoring drag. Which has greater speed at landing? (Hint: skip trajectories. Initial kinetic energy, final height difference, and gravitational transfer match. The upward arrow briefly stops at the top and takes a detour, but conservative forces recognize only endpoints.)
\item \textbf{A bullet and half a board.} A bullet just penetrates a board and exits with zero speed. Replace it with half the thickness of identical wood and fire the same bullet at the same speed. What speed remains? (Hint: half thickness means half the frictional work and half the kinetic loss. Half the energy remains, not half the speed; the speed fraction is \(1/\sqrt{2}\), about seventy percent. This answers difficulty two in the main text.)
\item \textbf{Pumped-storage accounting.} A plant uses surplus nighttime electricity to pump water from a lower to an upper reservoir, then releases it to generate at daytime peaks. The reservoirs differ by about \(500\) meters and store \(10^{10}\) kilograms of water. What theoretical electricity output, in kilowatt-hours, comes from emptying it? Why build a cycle with losses at every step? (Hint: \(E=mgh=10^{10}\times10\times500=5\times10^{13}\) joules divided by \(3.6\times10^{6}\) joules per kilowatt-hour gives about fourteen million kilowatt-hours. The losses buy time: unwanted midnight electricity becomes desirable daytime supply, perhaps losing twenty percent of energy while multiplying value. Conservation and economics keep different ledgers.)
\end{puzzles}

Next chapter brings another ledger. In collisions energy accounting can appear unbalanced: two lumps of modeling clay collide and much kinetic energy seems to disappear, yet another quantity remains unchanged. That quantity is momentum. Energy and momentum occupy two pages of the universe's ledger, neither replacing the other.
